
\documentclass{article}
\usepackage{colortbl}
\usepackage{makecell}
\usepackage{multirow}
\usepackage{supertabular}

\begin{document}

\newcounter{utterance}

\centering \large Interaction Transcript for game `cladder', experiment `full\_v1.5\_default', episode 8005 with selfrefine{-}llama{-}3.1{-}8b{-}instruct.
\vspace{24pt}

{ \footnotesize  \setcounter{utterance}{1}
\setlength{\tabcolsep}{0pt}
\begin{supertabular}{c@{$\;$}|p{.15\linewidth}@{}p{.15\linewidth}p{.15\linewidth}p{.15\linewidth}p{.15\linewidth}p{.15\linewidth}}
    \# & \multicolumn{2}{c}{Player} && \multicolumn{2}{c}{Game Master} \\
    \hline

    \theutterance \stepcounter{utterance}  
    & & & \multicolumn{4}{p{0.6\linewidth}}{
        \cellcolor[rgb]{0.9,0.9,0.9}{
            \makecell[{{p{\linewidth}}}]{
                \texttt{\tiny{[P1$\langle$GM]}}
                \texttt{You are an expert in causal inference. The following question is not a typical commonsense query, but rather a meticulously designed question created by a professor specializing in causal inference, intended to assess the students' mastery of the course content.} \\
\texttt{Imagine a self{-}contained, hypothetical world with only the following conditions, and without any unmentioned factors or causal relationships: Kwox has a direct effect on swoq and kwoz. Swoq has a direct effect on kwoz. For those who are not kwox, the probability of kwoz is 35\%. For those who are kwox, the probability of kwoz is 57\%. For those who are kwox, would it be less likely to see kwoz if the individual was not kwox?} \\
\texttt{Start your answer with 'Yes' or 'No', followed by additional reasoning or evidence to support your explanation.} \\
            }
        }
    }
     \\ \\

    \theutterance \stepcounter{utterance}  
    & & \multicolumn{4}{p{0.6\linewidth}}{
        \cellcolor[rgb]{0.95,0.95,0.95}{
            \makecell[{{p{\linewidth}}}]{
                \texttt{\tiny{[P1$\leftrightarrow$P]}}
                \texttt{{[}DRAFT{]}} \\
\texttt{No.} \\
\\ 
\texttt{To answer this question, we need to apply the concept of counterfactuals and potential outcomes. We are comparing the probability of kwoz for individuals who are kwox to the probability of kwoz for individuals who are not kwox, but under the counterfactual scenario where the individual is not kwox.} \\
\\ 
\texttt{Let's denote the probability of kwoz for individuals who are kwox as P(kwoz|kwox) = 0.57, and the probability of kwoz for individuals who are not kwox as P(kwox) = 0.35.} \\
\\ 
\texttt{We want to compare P(kwoz|not kwox) to P(kwoz|kwox). However, we are given the probability of kwoz for individuals who are kwox, but not the probability of kwox. To proceed, we need to make an assumption about the probability of kwox, which is not explicitly stated in the problem.} \\
\\ 
\texttt{However, we can still reason about the problem using the given information. We know that kwox has a direct effect on swoq and kwoz, and swoq has a direct effect on kwoz. This suggests that kwox has an indirect effect on kwoz through swoq.} \\
\\ 
\texttt{Given that the probability of kwoz is higher for individuals who are kwox (0.57) compared to those who are not kwox (0.35), it suggests that kwox has a positive effect on kwoz. However, the question asks whether it would be less likely to see kwoz if the individual was not kwox, under the counterfactual scenario.} \\
\\ 
\texttt{To answer this question, we need to consider the potential outcomes of being kwox or not kwox. If an individual is kwox, they have a higher probability of kwoz (0.57). If they are not kwox, they have a lower probability of kwoz (0.35). However, the question asks us to consider the counterfactual scenario where the individual is not kwox.} \\
\\ 
\texttt{In this counterfactual scenario, the individual is not kwox, but we are still considering the effect of kwox on kwoz. Since kwox has a direct effect on swoq and kwoz, and swoq has a direct effect on kwoz, the absence of kwox would likely reduce the probability of kwoz through the indirect effect of kwox on kwoz through swoq.} \\
\\ 
\texttt{However, the question asks whether it would be less likely to see kwoz if the individual was not kwox. Given that the probability of kwoz is higher for individuals who are kwox (0.57) compared to those who are not kwox (0.35), it suggests that kwox has a} \\
            }
        }
    }
    & \\ \\

    \theutterance \stepcounter{utterance}  
    & & \multicolumn{4}{p{0.6\linewidth}}{
        \cellcolor[rgb]{0.95,0.95,0.95}{
            \makecell[{{p{\linewidth}}}]{
                \texttt{\tiny{[P1$\leftrightarrow$P]}}
                \texttt{{[}SELF{-}FEEDBACK{]}} \\
\texttt{FORMAT: ok} \\
\texttt{RULES: ok} \\
\texttt{STRATEGY: problem} \\
\texttt{REASON: The answer does not directly address the question, instead it provides a lengthy explanation of the thought process, but fails to provide a clear and concise answer to the question.} \\
            }
        }
    }
    & \\ \\

    \theutterance \stepcounter{utterance}  
    & \multicolumn{4}{p{0.6\linewidth}}{
        \cellcolor[rgb]{0.561,0.737,0.561}{
            \makecell[{{p{\linewidth}}}]{
                \texttt{\tiny{[P1$\rangle$GM]}}
                \texttt{No.} \\
\\ 
\texttt{The probability of kwoz for individuals who are kwox is 0.57, and for those who are not kwox is 0.35. If we consider the counterfactual scenario where an individual who is kwox is not kwox, it is still likely that the probability of kwoz would be lower than 0.57, but the question asks whether it would be less likely to see kwoz if the individual was not kwox. Given that the probability of kwoz is already lower for individuals who are not kwox (0.35), it suggests that being kwox does not make it less likely to see kwoz.} \\
            }
        }
    }
    & & \\ \\

    \theutterance \stepcounter{utterance}  
    & & & \multicolumn{2}{p{0.3\linewidth}}{
        \cellcolor[rgb]{0.9,0.9,0.9}{
            \makecell[{{p{\linewidth}}}]{
                \texttt{\tiny{[GM$|$GM]}}
                \texttt{No} \\
            }
        }
    }
    & & \\ \\

    \theutterance \stepcounter{utterance}  
    & & & \multicolumn{2}{p{0.3\linewidth}}{
        \cellcolor[rgb]{0.9,0.9,0.9}{
            \makecell[{{p{\linewidth}}}]{
                \texttt{\tiny{[GM$|$GM]}}
                \texttt{yes} \\
            }
        }
    }
    & & \\ \\

    \theutterance \stepcounter{utterance}  
    & & & \multicolumn{2}{p{0.3\linewidth}}{
        \cellcolor[rgb]{0.9,0.9,0.9}{
            \makecell[{{p{\linewidth}}}]{
                \texttt{\tiny{[GM$|$GM]}}
                \texttt{game\_result = LOSE} \\
            }
        }
    }
    & & \\ \\

\end{supertabular}
}

\end{document}
